\section{Equivariance, differentiability, and design gradients}
\label{sec:differentiability}

A useful surrogate should support gradient-based geometry and circuit design.
Differentiability must therefore be considered at the level of the continuous
scene, not added after discretization.

\subsection{Rigid-motion symmetry}

For a common rigid transform $g\in\SE(3)$ in an isotropic environment,
\[
 Z(g\cdot\mathcal S)=Z(\mathcal S).
\]
For a continuous vector field,
\[
 E_{g\cdot\mathcal S}(g\mathbf r)
 =
 R_gE_{\mathcal S}(\mathbf r).
\]
The neural architecture enforces these transformations by relative poses,
local frames, and equivariant feature types rather than data augmentation
alone.

\subsection{Permutation symmetry}

Objects of the same semantic class have no arbitrary ordering. Scene encoders
must be permutation equivariant before global pooling. Port ordering is
physical, however, and transforms the port matrices by the corresponding
permutation:
\[
 Z'\!=P ZP^\T.
\]

\subsection{Geometry derivatives}

Let $\theta$ denote continuous shape/pose/material parameters. For a linear
integral equation
\[
 A(\theta)x(\theta)=b(\theta),
\]
the state derivative satisfies
\[
 A\frac{\pd x}{\pd\theta}
 =
 \frac{\pd b}{\pd\theta}
 -\frac{\pd A}{\pd\theta}x.
\]
For a scalar objective $J(x,\theta)$, an adjoint $\lambda$ solving
\[
 A\adj\lambda
 =
 \left(\frac{\pd J}{\pd x}\right)\adj
\]
gives
\[
 \frac{\dd J}{\dd\theta}
 =
 \frac{\pd J}{\pd\theta}
 +\lambda\adj
 \left(
  \frac{\pd b}{\pd\theta}
  -\frac{\pd A}{\pd\theta}x
 \right).
\]
Thus optimization gradients need not differentiate through every Krylov
iteration.

\subsection{Shape differentiability}

Superellipse and superquadric parameters are differentiable away from
degenerate limits. The geometry kernel exposes derivatives of coordinates,
normals, Jacobians, and pair distances. Collision/contact topology changes are
piecewise smooth events; optimization either enforces clearance constraints or
uses a smooth barrier before topology changes occur.

\subsection{Neural correction gradients}

FAST-mode gradients are obtained directly by automatic differentiation through
the scene encoder and structured decoders. CERTIFIED-mode gradients may combine
the neural derivative with an adjoint correction of the physical residual.
The correction is essential when optimization approaches the edge of the
training distribution.

\subsection{Gradient certification}

A design gradient is optionally checked by:
\begin{enumerate}
\item adjoint residual;
\item one or more random directional finite-difference probes;
\item symmetry identities under rigid scene perturbations.
\end{enumerate}
A low prediction residual does not automatically certify a derivative.
